\documentclass[5p]{elsarticle}

\usepackage{amsmath,amssymb}
\usepackage{booktabs}
\usepackage{graphicx}
\usepackage{url}
\usepackage[colorlinks=true,linkcolor=blue,citecolor=blue,urlcolor=blue]{hyperref}
\usepackage{xurl}
\usepackage{xcolor}

% ToUnicode maps, so that ligatures and dashes survive text extraction. Without
% these the submission system reads "Syafiq" as "Syaq" and loses the dash in
% "Writing -- review".
\usepackage{cmap}
\usepackage[T1]{fontenc}
\usepackage{lmodern}

\newcommand{\todo}[1]{\textcolor{red}{[TODO: #1]}}
\newcommand{\rhoSym}{\rho}
% Elsevier's two-column 5p layout is narrower than IEEEtran's, so the wide tables
% overflow by a few points. Shrink the table floats and column padding here rather
% than editing the shared body, which must stay identical to the IEEE version.
\setlength{\tabcolsep}{4pt}
\setlength{\emergencystretch}{3em}
\usepackage{etoolbox}
\AtBeginEnvironment{thebibliography}{\raggedright}
\AtBeginEnvironment{table}{\small}
\AtBeginEnvironment{tabular}{\small}

\journal{Journal of Systems Architecture}

\begin{document}

\begin{frontmatter}

\title{The cache is cheaper than the weights: measured bit costs and device-dependent capacity for LLM inference on 8\,GB devices}

\author[uum]{Yibo Zhang\corref{cor1}}
\ead{zhang\_yibo@ahsgs.uum.edu.my}
\cortext[cor1]{Corresponding author.}
\author[uum]{Muhammad Syafiq bin Mohd Pozi}
\author[uum]{Osman B Ghazali}

\affiliation[uum]{organization={School of Computing, Universiti Utara Malaysia},
            city={Sintok},
            state={Kedah},
            country={Malaysia}}

\begin{abstract}
Quantisation fits a large language model into an 8\,GB device, and the
literature optimises it for speed. We measure four mismatches between that
literature and the device. \emph{First}, the serving stack sets the budget before
any bit does: grouped-query attention silently disables every fused attention
kernel on our stack, the fallback materialises an $\mathrm{O}(L^2)$ score matrix,
and materialising the KV heads --- a mathematically transparent rewrite --- cuts
attention peak memory $32\times$ and 4k prefill time $5.8\times$. \emph{Second},
4-bit weights buy little or no time, because dequantisation is not fused into the
matmul; quantisation here buys capacity, not throughput, so the objective must be
capacity. \emph{Third}, the per-bit cost runs against practice: with the published
GPTQ for weights and KIVI's own quantiser for the cache, weights cost $11.39\%$
perplexity where the cache costs $1.27\%$ on Qwen3-0.6B, and $3.20\%$ against
$0.24\%$ on Qwen3-1.7B; the ordering reproduces on WikiText-2, and a 2-bit passkey
experiment shows the scheme, not the width alone, decides whether long-context
retrieval survives. \emph{Fourth}, capacity is device-dependent: usable memory
differs $3.15\times$ between a discrete-GPU laptop and a unified-memory edge
module. We solve the allocation exactly and report both devices' allocations and
the module's measured ceiling.
\end{abstract}

\begin{keyword}
quantisation \sep LLM inference \sep KV cache \sep edge inference \sep
memory budget \sep bit allocation \sep energy efficiency
\end{keyword}

\end{frontmatter}

% --------------------------------------------------------------------------
% Highlights, as Elsevier requires (3-5 bullets, 85 characters each at most).
\section*{Highlights}
\begin{itemize}
\item On an 8\,GB device the serving stack, not the bit width, sets the memory budget.
\item One KV-cache bit costs 9--13x less accuracy than one weight bit, on two models.
\item 2-bit cache quantisation separates schemes where 4-bit does not: 5/5 against 0/5.
\item The measured ceiling of the edge module is 3k tokens, not the 67k a model predicts.
\item The optimal policy compresses the cache and leaves the weights at 8 bits.
\end{itemize}

\section*{CRediT authorship contribution statement}
\textbf{Yibo Zhang:} Conceptualization, Methodology, Software, Validation,
Formal analysis, Investigation, Data curation, Writing - original draft,
Writing - review \& editing, Visualization. \textbf{Muhammad Syafiq bin Mohd
Pozi:} Supervision. \textbf{Osman B Ghazali:}
Supervision.

\section*{Declaration of competing interest}
The authors declare that they have no known competing financial interests or
personal relationships that could have appeared to influence the work reported in
this paper.

\section*{Data availability}
All measurement scripts, the vendored third-party quantisers, the measurement
files behind every table, and the scripts that regenerate the figures are
publicly available at
\url{https://github.com/BRAD-cording/llm-8gb-bit-costs}. A single command rebuilds
the results and runs the checks described in Section~\ref{sec:repro}.

\section{Introduction}
\label{sec:intro}

Deploying a large language model on an 8\,GB device is no longer a research
curiosity; it is the default target for on-device assistants, and the entire
post-training quantisation (PTQ) literature is organised around it.

\emph{What kind of paper this is.} It is a measurement paper. Its subject is what
an 8\,GB device actually does with the bits it is given, and its contributions are
four measurements that contradict what the allocation literature assumes, plus the
allocator that follows from them. We state that plainly because the opposite
reading is available: the allocation step is a multiple-choice knapsack solved by
exact dynamic programming, which is textbook. Nothing here should be read as a new
optimiser. What is new is the input to the optimiser --- per-device usable memory,
a measured workspace curve, and per-bit accuracy costs measured on both sides of
the same model --- and the finding that the objective has to be capacity rather
than latency. Section~\ref{sec:method} returns to this distinction. That
literature optimises \emph{weights}: rotations to flatten activations
\cite{quarot,spinquant,duquant,sun2024}, Hessian-aware rounding
\cite{gptq,frantar2022}, and mixed-precision schemes that spend more bits on
sensitive layers \cite{hawqv2,mixllm,gamma,alignquant,wang2018}. System-level work
co-designs the kernel with the format \cite{lin2024,zhao2023,wei2024,guo2026,mixquant}, and
low-rank compensation has been used to recover what rounding loses
\cite{song2026,park2026,liu2025}. The KV cache is either left in 16 bits or
compressed by a separate line of work \cite{kivi,kvquant,mixkvq,kvpareto,duanmu2024,
su2025,he2024} that takes the weight precision as given. Serving stacks
that manage the cache as a first-class resource \cite{kwon2023,kwon2023} do not
choose weights against cache either.

At the context lengths that make these models useful, that division is the wrong
one. From the published model cards this is arithmetic rather than a hazard: a KV cache costs 112\,KiB per token for Qwen3-0.6B and 1.7B and 144\,KiB for 4B and 8B, so at 32k tokens the cache is 3.76--4.83\,GB against 0.22--4.09\,GB of 4-bit weights, and the two are equal at 1{,}919 tokens for 0.6B, 6{,}144 for 1.7B, 12{,}319 for 4B and 21{,}388 for 8B. For an 8\,GB device the cache overtakes the weights long before the context lengths these models are deployed at.

Once the cache is a first-class consumer, the problem becomes an allocation
problem: given a fixed byte budget, how many bits for weights and how many for
the cache? This paper argues that the answer is \emph{device-dependent}, and
that the dependence has a cause that can be measured but that no PTQ paper
reports.

The cause is the gap between nominal and usable memory. Both of our devices are
advertised as 8\,GB. Measured on the day of the experiment, a discrete-GPU
laptop has 6.93\,GiB free for the model ($\rhoSym=0.133$: driver context,
compositor) while a unified-memory Jetson Orin Nano Super has 4.32\,GiB
($\rhoSym=0.419$: L4T, desktop, carve-out). The usable budgets differ by
$1.6\times$. Because the cache grows linearly in context length while the
weights do not, that difference changes the exchange rate between a weight bit
and a token of context: dropping Qwen3-4B from 8-bit to 4-bit weights multiplies
the attainable context by $1.57\times$ on the laptop but by $5.76\times$ on the
Jetson --- a $3.7\times$ difference in the price of the same bit
(Section~\ref{sec:rho}). A single published configuration cannot be optimal on
both.

\paragraph*{Contributions}
\begin{enumerate}
\item \textbf{What sets the budget is measurable, and it is not the bit-width.} A
      silent attention-dispatch failure costs $32\times$ in attention peak memory
      and $5.8\times$ in 4k prefill time, and removing it extends the attainable
      context $4\times$ before any quantisation decision is made; separately, the
      usable-memory fraction differs $3.15\times$ between two nominally identical
      8\,GB targets. Both are properties of the stack and the device that the
      allocation literature treats as fixed. We release a model-free reproduction
      of the dispatch failure, which leaves no trace in the framework's output, and
      we report the edge module's \emph{measured} ceiling beside the cost model's
      prediction for it (Sections~\ref{sec:results-stack} and~\ref{sec:rho}).

\item \textbf{The per-bit cost of the two sides runs against practice.} With
      published implementations on both sides --- the official GPTQ for weights,
      KIVI's own quantiser for the cache --- one cache bit costs far less accuracy
      than one weight bit: $9.0\times$ at 0.6B and $13.3\times$ at 1.7B, and the
      ordering reproduces on WikiText-2. A 2-bit passkey experiment locates where
      the claim stops: at 4 bits the quantisation scheme does not change
      long-context retrieval, at 2 bits it decides it (Section~\ref{sec:results-bitcost}).

\item \textbf{The consequence is an allocation problem whose objective must be
      capacity.} We formulate joint weight--KV allocation as a multiple-choice
      knapsack over measured per-device curves and solve it exactly, so the
      optimisation is textbook --- the contribution is the inputs it runs on and
      the finding that latency cannot be the objective on this stack. The
      resulting policy compresses the cache and leaves the weights at 8 bits on
      both devices, and we report it against a deployed control rather than a
      modelled one (Sections~\ref{sec:method}--\ref{sec:results-capacity}).
\end{enumerate}

\paragraph*{What this paper does not claim}
We do not claim a new quantiser: we use existing ones and allocate their output.
We do not claim that the optimal allocation reverses between prefill and decode
regimes --- we tested that hypothesis and could not support it; the reasoning and
the falsification are documented in the artifact
(Appendix~\ref{app:prereg}). And we do not claim exhaustive literature coverage.

% --------------------------------------------------------------------------- %
\section{Related Work}
\label{sec:related}

\subsection{Quantisation for speed}
The dominant framing in post-training quantisation is that fewer bits mean
faster inference. Rotations flatten activations so a single 4-bit quantiser
suffices \cite{quarot,spinquant,duquant,prefixquant,ostquant,spectrumquant};
second-order methods cancel rounding error \cite{gptq,gptaq}; system papers
report speedups from weight-only formats \cite{awq}. That framing is what our
Section~\ref{sec:results-nospeed} measures directly, and on our stack it does
not hold: the arithmetic stays fp16, so the only quantity that changes is bytes
stored.

\subsection{Capacity-oriented compression}
A second line treats memory, not time, as the scarce resource: KV cache
quantisation \cite{kivi,kvquant,mixkvq}, joint memory--accuracy frontiers
\cite{kvpareto}, and precision allocation under a storage budget
\cite{hawqv2,mixllm,gamma}. Closest to our setting, \cite{alignquant} scores
precision choices separately for prefill and decode and \cite{dq} specialises
formats per phase, while \cite{mixquant} keeps decode in high precision.

Reference~\cite{kvpareto} deserves a precise comparison because it is the nearest prior
art. It maps a joint frontier over cache bit-width, cache granularity, chunked
prefill and weight quantisation, and reports Pareto-optimal configurations at
68--78\% memory reduction. Its weight side is fixed at 4-bit AWQ, however, so the
weight/cache \emph{split} is not among its decision variables; its frontier is
indexed by memory consumed rather than by a per-bit accuracy cost measured end to
end; and it reports no usable-memory fraction for its hardware, so the frontier
cannot be re-targeted to a device with a different $\rhoSym$. Our claim is
narrower and different: that the per-bit cost of the two sides must be measured
on the same footing, that the resulting optimum is device-dependent through
$\rhoSym$, and that on our stack the answer is to compress the cache and leave
the weights alone.

None of these treats the weight/KV split as a decision variable under a measured
device budget, and none reports the usable-memory fraction of the target device,
which is the quantity that dominates our results.

\subsection{Edge measurement and reporting}
A recurring criticism is that quantisation results are not reproducible because
the execution stack is unreported. Our measurements give that criticism a
number: the same model on the same card supports 4k or beyond 16k context
depending only on which attention kernel the framework selects
(Section~\ref{sec:results-stack}). This is also why we record the software stack
and the device's memory state with every latency number, and why our $\rhoSym$
values are reported as a measurement procedure rather than as a constant. What we
do \emph{not} control is die temperature and clock throttling: our sweeps are
short and configurations are interleaved, which bounds thermal ordering effects,
but we did not instrument the clocks, so a throttling explanation for a small
difference cannot be excluded from our data alone.

% --------------------------------------------------------------------------- %
\section{Usable Memory and the Weight-Bit Exchange Rate}
\label{sec:rho}

\paragraph*{What ``8\,GB'' actually means}
We measure the fraction of nominal memory that is unavailable to the model on
the day of the experiment, rather than trusting the datasheet. On the laptop this
is the CUDA context plus the compositor, read from the allocator; on the Jetson it
is L4T, the desktop and the carve-out, read from the kernel. The difference is
large enough to change conclusions:



The usable budgets differ by $1.6\times$ although both devices are sold as
8\,GB. $\rhoSym$ is not a constant: it moves with the OS image, with background
load, and --- on the Jetson --- with the power mode, which we set once
(25\,W) and did not vary, so the power-mode sensitivity is a stated caveat
rather than a measurement. This is why we report the procedure alongside the
number. Any paper that states a device and a bit-width
without stating $\rhoSym$ is reporting a configuration that does not transfer.

\paragraph*{The exchange rate between a weight bit and a token}
Because the cache grows linearly in context and the weights do not, a saved
weight bit converts into context at a rate that depends on how tight the budget
is. Figure~\ref{fig:exchange} shows the conversion for Qwen3-4B with the cache
held at 16 bits: dropping weights from 8 to 4 bits multiplies the attainable
context by $1.57\times$ on the laptop and by $5.76\times$ on the Jetson, so the
same bit is worth $3.7\times$ more on the constrained device. The rates are
computed from the measured $\rhoSym$ and workspace curves of
Section~\ref{sec:method}, not from datasheet numbers.

\begin{figure}[t]
\centering
\includegraphics[width=0.8\columnwidth]{figures/fig_exchange.pdf}
\caption{Attainable context versus weight width for Qwen3-4B with a 16-bit cache,
from measured $\rhoSym$ and workspace. The two curves have different slopes,
which is the whole reason a single published configuration cannot be optimal on
both devices.}
\label{fig:exchange}
\end{figure}

\paragraph*{Why this is a capacity argument, not a speed argument}
It would be natural to read the curves above as ``quantise the weights to go
faster''. They do not say that: Section~\ref{sec:results-nospeed} shows that the
4-bit path is in fact slower on this stack. What the curves say is that on a
device this constrained, precision is the knob that decides how much context you
can address at all, and that the exchange rate for that knob is device-dependent.
That is the sense in which the allocation problem is worth solving.

% --------------------------------------------------------------------------- %
\section{Method}
\label{sec:method}

\paragraph*{Why the objective is capacity, not latency}
The natural formulation is to minimise serving time subject to an accuracy
constraint. We tried that first and it is degenerate on this stack: latency is
not a function of the bit-widths at all, since a 4-bit weight is still multiplied
at fp16 precision (Table~\ref{tab:nospeed}). An objective whose gradient with
respect to the decision variables is zero cannot rank allocations. What bit-width
\emph{does} determine here is how much context fits, so that is what we optimise:
maximise the attainable context subject to a bounded accuracy loss. This is a
measured choice, not a modelling convenience, and it is the reason
Table~\ref{tab:nospeed} appears before the method rather than after it.

\paragraph*{Budget constraint}
Let $M$ be nominal memory and $\rhoSym$ the measured unusable fraction. An
allocation $\mathbf b$ is feasible for a context length $L$ iff
\begin{equation}
\sum_{l} \Big( b^{W}_l \frac{|W_l|}{8} + s_l \Big)
+ 2 L d H_{kv} \frac{\bar b^{KV}}{8}
\;\le\; M(1-\rhoSym) - W_{\mathrm{act}}(L),
\label{eq:budget}
\end{equation}
where $s_l$ are per-group scales and the cache term is linear in $L$. The
workspace term is the one prior budget model gets wrong: it is neither zero nor
constant. Measured, it is linear in the tokens of a single forward pass with a
\emph{model-dependent} coefficient and intercept: $0.142$\,MB per token and
9\,MB for the 0.6B model, $0.217$ and 161\,MB for the 4B. The coefficient has to
be refitted per model --- using the 4B value for a 0.6B model over-reserves
workspace by about $60\%$ and so under-reports attainable context, a mistake we
made and corrected. What it does not depend on is the device: the 0.6B curve
fitted on the laptop predicts the Jetson's measured workspace to within $2\%$
(Section~\ref{sec:setup}). Because the term scales with the tokens in a single
forward pass rather than with the conversation, chunked prefill bounds it at
$W_{\mathrm{act}}(L) = c \cdot \min(L, C) + c_0$ with chunk size $C$. Fitting the
whole prompt in one pass is what makes 32k look infeasible on an 8\,GB card; it
is a property of the serving stack, not of the device.

\paragraph*{Objective}
For a target context length $L$,
\begin{equation}
\min_{\mathbf b} \; \sum_i \Delta\mathrm{PPL}_i(b_i)
\quad \text{s.t.} \quad \eqref{eq:budget},
\label{eq:obj}
\end{equation}
where the sum runs over every weight matrix and every KV group and
$\Delta\mathrm{PPL}_i(b)$ is the perplexity cost of quantising unit $i$ to $b$
bits, measured per unit (Section~\ref{sec:setup}). The capacity question
\emph{at an accuracy budget} $\epsilon$ is the inverse: the largest $L$ for which
\eqref{eq:obj} attains a value below $\epsilon$. We report both directions, since
the second is what a deployment engineer asks and the first is what shows
\emph{where} the bits went.

\paragraph*{Solver}
With $\mathbf b$ discrete per unit, \eqref{eq:obj} is a multiple-choice knapsack
and dynamic programming over discretised byte buckets returns the exact optimum
in $O(n \cdot |\mathcal{B}| \cdot \text{buckets})$; no heuristic search is
needed, and the full Pareto front comes out in one pass. Two implementation
details matter for correctness and cost us two bugs to learn: bucket costs must
be rounded \emph{up} (rounding down lets the solver quietly exceed the budget),
and the bucket grid must be finer than the number of decision units (otherwise
every unit's minimum cost of one bucket makes a feasible problem look
infeasible).

\paragraph*{What is measured and what is assumed}
$\rhoSym$, $W_{\mathrm{act}}(L)$, $B_{\mathrm{eff}}$ and $P_{\mathrm{eff}}$ are
measured per device. The per-unit costs $\Delta\mathrm{PPL}_i$ are measured for
every linear layer at 2, 3 and 4 bits, and for the cache globally; because the
per-unit measurement is noise-limited for most layers at 4 bits, we keep the
measured \emph{ranking} and rescale so that the sum matches the end-to-end
measured curve. Spreading the cache cost uniformly across layers is an
assumption, stated as one, because a per-layer cache measurement was not
available.

% --------------------------------------------------------------------------- %
\section{Experimental Setup}
\label{sec:setup}

\paragraph*{Devices}
Two nominally 8\,GB targets, chosen because they bracket the two memory
architectures an edge deployment realistically uses.
\emph{Device A} is an RTX 4060 Laptop GPU (8\,188\,MiB, sm\_89, driver 592.82,
measured 161.9\,GB/s copy bandwidth and 22.36\,TFLOP/s FP16 at batch 128).
\emph{Device B} is a Jetson Orin Nano Super (7\,619\,MiB unified LPDDR5,
L4T R36.4.7, sm\_87, 8 SMs, measured 54.1\,GB/s copy bandwidth and
4.51\,TFLOP/s FP16 at batch 128). We report its \emph{usable memory} as measured --- 7\,619\,MiB unified against
4.32\,GiB actually usable ($\rhoSym{=}0.419$; the laptop is 8\,188\,MiB against
6.93\,GiB, $\rhoSym{=}0.133$) --- and we state plainly which of its numbers are measured and
which are derived, because the distinction matters:

\begin{itemize}
\item \textbf{Measured on Device B}: total and available memory, hence
      $\rhoSym$; copy bandwidth ($54.1$\,GB/s) and achieved FP16 throughput
      ($4.51$\,TFLOP/s at batch 128), hence $b^*_{\mathrm{Jetson}}{=}83$; and an
      end-to-end sweep of the real 0.6B checkpoint at 512--2{,}048 tokens
      (TTFT 165.1, 307.6 and 638.0\,ms; $\sim$108\,ms per decode step; peaks of
      1{,}334, 1{,}468 and 1{,}734\,MB).
\item \textbf{Derived, not measured, on Device B}: the context lengths in its
      column of Table~\ref{tab:capacity}, and its curve in
      Figure~\ref{fig:exchange}. Those follow from the measured $\rhoSym$
      together with a workspace curve fitted on Device A. That transfer is now
      evidence-backed rather than assumed: the fitted curve $0.142L+9$\,MB
      predicts the three measured Jetson peaks to within $2\%$ (83, 158 and
      307\,MB against predictions of 82, 154 and 300), so the workspace is a
      property of the model and the framework rather than of the GPU. The
      \emph{lengths} remain analytic because the naive attention path exceeds
      the module's usable memory above 2k tokens, which is itself the
      Section~\ref{sec:results-stack} effect.
\end{itemize}

We would rather ship a paper that says which of its numbers are measured than one
whose tables imply a second set of measurements that do not exist.

\paragraph*{Software stack}
Every number carries the stack it was measured on, because
Section~\ref{sec:results-stack} shows that the stack can move attainable context
by $4\times$. Unless stated otherwise: torch 2.6.0+cu124, transformers 5.20.2,
bitsandbytes 0.50.2, Windows 11 on Device A and L4T on Device B. Weights come
from ModelScope; the network in our environment cannot reach HuggingFace, and
\S\ref{sec:setup} of the artifact records the package indices used.

\paragraph*{Measurement protocol}
Each configuration is measured by (i) loading and warming up, (ii) timing three
prefill passes over the full prompt and taking the median, (iii) timing eight
decode steps against a pre-built cache and taking the median. Timings are taken
under \texttt{torch.inference\_mode}; prefill is timed on the base model so that
the language-model head is not applied to every position
(App.~\ref{app:trap}). Peak memory is read from the CUDA allocator after each
phase separately, and any run whose peak exceeds $0.92$ of device memory is
flagged as spilled and excluded rather than reported: on Windows the driver backs
an over-budget allocation with host memory instead of raising an error, which
degrades throughput by up to $30\times$ silently.

\paragraph*{Interleaving}
Configurations that are compared are measured \emph{interleaved}, not one after
the other. Background load moved a short-context prefill time by $10\%$ between
two consecutive runs of the same configuration in our own measurements; an
earlier draft compared nf4 and fp16 measured minutes apart and reported the
wrong ratio. All ratios in Table~\ref{tab:nospeed} are from interleaved runs.

\paragraph*{Models}
Qwen3-0.6B, Qwen3-1.7B and Qwen3-4B, with configurations taken from the
published model cards, whose per-token cache sizes are 112\,KiB (0.6B, 1.7B) and 144\,KiB (4B, 8B). The 0.6B and 1.7B models carry
the accuracy measurements, where a full-precision reference still fits; the 4B
model carries the system measurements, where one does not. Sizes above 4B are
analysed but not run: their 4-bit weights plus a 16-bit cache exceed 8\,GB
before any context is added, which is itself the motivation for the allocation
problem.

\paragraph*{What is synthetic, and why it is legitimate}
Once the real 0.6B checkpoint arrived we found that an architecture-exact
random-weight model reproduces its memory figures within $1\%$ and its naive
prefill time within $4\%$, because latency and footprint depend on tensor shapes
and dtypes rather than on weight values. The synthetic path let us run the stack
ablation before a 1.5\,GB download completed on a $0.2$\,MB/s link. All numbers
in Tables~\ref{tab:stack} and~\ref{tab:nospeed} are from real weights; synthetic
values are kept in the artifact as a cross-check.

\paragraph*{Accuracy measurement}
Perplexity is measured on a 575k-token corpus assembled from 2{,}000 arXiv
cs.CL/cs.LG abstracts, and layer-wise sensitivity by quantising one linear
layer at a time at 2, 3 and 4 bits (591 measurements, 197 layers). The corpus
is a documented fallback: WikiText-2 could not be obtained because no Western
host is reachable from this network, so absolute perplexities are not
comparable to the literature and only relative increases are used. The
sensitivity measurement is noise-limited for most layers at 4 bits (the floor is
$0.12\%$ relative), which is why the allocator keeps the measured ranking and
rescales it to the end-to-end curve.

% --------------------------------------------------------------------------- %
\section{Results}
\label{sec:results}

Every system number below is measured on the target hardware with a real
checkpoint. Table~\ref{tab:capacity} is the exception and is analytic by
construction: it follows from the measured inputs named in
Section~\ref{sec:setup}, and its workspace model is validated on both devices.
Accuracy numbers use the abstract-level corpus and are reported as relative
perplexity only.

\subsection{The stack decides the budget before any bit does}
\label{sec:results-stack}

On this stack --- RTX 4060 Laptop, Windows, torch 2.6.0+cu124, transformers
5.20.2 --- a Qwen3-style model never reaches a fused attention kernel. Its 16
query heads and 8 KV heads violate the equal-head requirement of the flash and
memory-efficient backends, so PyTorch dispatches to the math backend, which
materialises the score matrix in fp32. The failure is silent: no exception, no
warning that survives into a log.

Table~\ref{tab:gqa} isolates the cause. With identical tensors and identical
mathematics, materialising the KV heads makes the fused kernel available and
reduces the attention peak by $32\times$; we verified numerically that the two
paths agree.

\begin{table}[t]
\centering
\caption{Attention peak memory, $q{=}16$ heads, $kv{=}8$ heads, $L{=}4096$,
$d{=}128$, fp16. Identical output values in every successful row.}
\label{tab:gqa}
\begin{tabular}{llr}
\toprule
KV heads & backend & peak \\
\midrule
as-is & memory-efficient & \emph{no kernel} \\
as-is & flash & \emph{no kernel} \\
materialised & memory-efficient & \textbf{84 MB} \\
materialised & math & \textbf{2{,}693 MB} \\
\bottomrule
\end{tabular}
\end{table}

Table~\ref{tab:stack} shows what this does to end-to-end serving. Restoring the
fused path moves the attainable context from 4k (with a silent spill at 8k) to
beyond 16k, and cuts prefill time at 4k by $5.8\times$. At 8k the naive path's
71\,s is a wall-clock figure for a spilled run and we do not treat it as a
latency measurement; the fused path completes the same prefill in 0.99\,s, a
$72\times$ wall-clock difference that a deployment would feel even though it is
not a clean latency ratio. At the 4B scale the effect is qualitative rather than
quantitative: with 4-bit weights and the fallback, a 4k prompt already exceeds
the card (9{,}195\,MB peak, measured with the real checkpoint) and the run
degrades to 47\,s per prefill, whereas the fused path completes the same prefill
in 1.9\,s with 2.7\,GB to spare. That
is not a slow configuration but an unusable one, and nothing in the framework
reports it.

The fix is not free everywhere, and the second device shows the boundary.
Repeating the sweep on the Jetson gives 165.1, 307.6 and 638.0\,ms of TTFT for
the naive path at 512, 1{,}024 and 2{,}048 tokens, against 189.4, 354.6 and
833.3\,ms with the KV heads materialised --- the rewrite is $15$--$31\%$
\emph{slower} there. At these lengths the $\mathrm{O}(L^2)$ term is too small to
pay for the Python-level repeat, and the module's CPU is slow enough that the
overhead shows. On the laptop the crossover falls between 2k and 4k, which is
also where the naive path starts to spill; on the Jetson we cannot reach 4k at
all, because the framework materialises logits for every position ($1.2$\,GB at
$L{=}4096$ for a 152k vocabulary) on top of a model and cache that already
occupy half the usable memory. Applying the dispatch fix below roughly 2k tokens
is therefore a pessimisation, and it should be gated on context length rather
than enabled globally --- a caveat we did not anticipate when we first measured
the $32\times$ memory reduction.

\begin{table}[t]
\centering
\caption{End-to-end effect of the attention dispatch, on real checkpoints at both
sizes. We first screened the 4B scale with an architecture-exact random-weight
model and then re-ran every entry with the real checkpoint: the synthetic model
predicted peak memory to within $0.7\%$ and prefill time to within $8\%$, so the
screening procedure is sound, but every number below is real-weight.
``spilled'' marks runs that exceeded VRAM and were backed by system memory
without raising an error, so their timings are not usable as latency.}
\label{tab:stack}
\begin{tabular}{rrrr}
\toprule
$L$ & TTFT naive & TTFT fused & peak naive \\
\midrule
\multicolumn{4}{c}{\emph{Qwen3-0.6B, fp16}}\\
1{,}024 & 164.7 ms & \textbf{62.2 ms} & 1{,}644 MB \\
4{,}096 & 2{,}062.5 ms & \textbf{355.9 ms} & 4{,}836 MB \\
8{,}192 & 71{,}420 ms (spilled) & \textbf{985.0 ms} & 13{,}436 MB \\
16{,}384 & out of memory & \textbf{2{,}795.2 ms} & --- \\
\midrule
\multicolumn{4}{c}{\emph{Qwen3-4B, nf4}}\\
1{,}024 & 692.3 ms & \textbf{425.3 ms} & 3{,}433 MB \\
4{,}096 & \emph{spilled} (46.7 s) & \textbf{1{,}892 ms} & 9{,}195 MB \\
8{,}192 & --- & \textbf{4{,}356 ms} & 5{,}795 MB \\
\bottomrule
\end{tabular}
\end{table}

\begin{figure}[t]
\centering
\includegraphics[width=\columnwidth]{figures/fig_stack.pdf}
\caption{Time to first token versus context length. Crosses mark runs that spilled or ran out of memory; they are drawn rather than dropped, because where the naive path stops working is the point.}
\label{fig:stack}
\end{figure}

\subsection{Quantisation buys capacity, not speed}
\label{sec:results-nospeed}

Table~\ref{tab:nospeed} compares 4-bit nf4 weights against fp16 on the real
checkpoint with the same attention path, measured in an interleaved fashion so
that both configurations see the same background load. nf4 is \emph{slower} at
every length, by $21\%$ at 1k and $7\%$ at 16k. The mechanism is that nf4 stores
4-bit weights but performs an fp16 tensor-core matmul, so the FLOPs are
identical, while the unfused dequantisation adds elementwise traffic that the
saved bytes do not pay for.

\begin{table}[t]
\centering
\caption{4-bit nf4 versus fp16 on the real Qwen3-0.6B checkpoint, fused
attention path. Every row is the mean of two interleaved runs, so both
configurations saw the same background load. nf4 is slower at every length, and
the penalty shrinks with context as the cache rather than the weights comes to
dominate traffic.}
\label{tab:nospeed}
\begin{tabular}{rrrrr}
\toprule
& \multicolumn{2}{c}{TTFT} & \multicolumn{2}{c}{TPOT} \\
\cmidrule(lr){2-3}\cmidrule(lr){4-5}
$L$ & fp16 & nf4 & fp16 & nf4 \\
\midrule
1{,}024 & 65.4 ms & 79.4 ms ($1.21\times$) & 26.9 ms & 37.1 ms \\
4{,}096 & 356.9 & 404.2 ($1.13\times$) & 28.8 & 37.2 \\
8{,}192 & 974.5 & 1{,}063.4 ($1.09\times$) & 50.7 & 71.5 \\
16{,}384 & 2{,}765.0 & 2{,}945.1 ($1.07\times$) & 75.5 & 75.4 \\
\bottomrule
\end{tabular}
\end{table}

A third precision makes the point sharper still. Bit-width does not even
\emph{order} latency on this stack: bnb's 8-bit path stores 752\,MB against
nf4's 531\,MB, yet costs 520.9\,ms of TTFT at 4k against nf4's 404.2\,ms and
fp16's 355.9\,ms, and 129.2\,ms per decode step against 37.2 and 28.8. The
reason is that bnb's int8 path performs per-multiply outlier decomposition,
which is a property of that kernel and not of 8-bit arithmetic.

Nor is the ordering stable across phases or model sizes. At 4B the same
comparison reverses for prefill and widens for decode: int8 (4.41\,GB of
weights) prefills 4k in 1{,}613\,ms against nf4's 1{,}744\,ms, while decoding at
128.5\,ms per token against 54.7. A single ``quantisation speedup'' number is
therefore not meaningful without naming the phase and the model size --- which is
precisely why an allocation objective written in terms of latency cannot be used
here.

The consequence for method design is direct: an allocation objective written in
terms of latency is degenerate on this stack, because latency does not depend on
the decision variables. We therefore optimise capacity subject to an accuracy
budget.

\subsection{Capacity is device-dependent}
\label{sec:results-capacity}

Table~\ref{tab:capacity} solves the allocation problem on both devices for a
range of accuracy budgets. At every budget the laptop reaches substantially more
context, by a factor that turns out to be almost constant (see below).

\begin{table}[t]
\centering
\caption{Largest context length whose minimum achievable perplexity increase
stays within $\epsilon$, computed from the measured $\rhoSym$ and workspace
curves and the measured weight/cache curves of Table~\ref{tab:bitcost}. The
Jetson column is an \emph{upper bound}, derived from that device's measured
$\rhoSym$ and a workspace curve validated on it to within $2\%$
(Section~\ref{sec:setup}). The extrapolation is not achieved in practice: direct
measurement puts the Jetson's ceiling for 0.6B near 3k tokens, and a 4k prefill
aborts the allocator outright. We report the column
because the comparison between devices is the point, and flag it because a reader
who took 67k as a working configuration would be misled. The allocator keeps weights at 8 bits
and spends the budget on the cache at every budget shown, because a cache bit is
not more expensive in accuracy and, at these context lengths, saves more bytes.}
\label{tab:capacity}
\begin{tabular}{rrrrrr}
\toprule
& \multicolumn{2}{c}{max context} & & \multicolumn{2}{c}{bits chosen} \\
\cmidrule(lr){2-3}\cmidrule(lr){5-6}
$\epsilon$ & 4060 & Jetson & ratio & weights & cache \\
\midrule
0.01 & 111{,}184 & 67{,}324 & $1.65\times$ & 8.00 & 7.7 / 7.4 \\
0.02 & 122{,}404 & 73{,}444 & $1.67\times$ & 8.00 & 6.9 / 6.6 \\
0.05 & 176{,}464 & 108{,}124 & $1.63\times$ & 8.00 & 4.8 / 4.6 \\
0.10 & 221{,}344 & 121{,}384 & $1.82\times$ & 7.98 / 8.00 & 3.8 / 4.1 \\
0.20 & 243{,}784 & 143{,}824 & $1.70\times$ & 7.98 / 8.00 & 3.3 / 3.4 \\
\bottomrule
\end{tabular}
\end{table}

\paragraph*{The Jetson column is an upper bound, not a recipe}
We checked the extrapolation directly and it does not survive. We ran a prefill of increasing length on the Jetson with the fused attention path
enabled: 1{,}024 tokens run at a peak of 1.52\,GB, 2{,}048 at 1.92\,GB and 3{,}072 at 2.32\,GB, while 4{,}096 aborts rather than degrading. The failure is not a
clean OOM exception --- it is an NVML assertion inside the caching allocator,
after the driver reports \texttt{NvMapMemAllocInternalTagged error 12} --- which is
why the bound can only be located by running each length in its own process.



This changes how the capacity result should be read, and we would rather say so
than let the number stand. The model behind Table~\ref{tab:capacity} contains the
measured $\rhoSym$ and a workspace curve fitted to that device, and it reproduces
the 512--2{,}048 range to within $2\%$. What it does not capture is that this
device fails \emph{in blocks}: the allocator aborts on a large allocation instead
of paging or fragmenting its way through, so the smooth extrapolation is simply
not the operative mechanism past a few thousand tokens. The laptop, which spills
instead of aborting, does run at 32k (Table~\ref{tab:headtohead}), which is the
same failure mode appearing benign on one device and fatal on the other.

The comparison between devices therefore stands as a statement about the cost
model, and the Jetson column should be read as the budget the model believes the
device has, not as a configuration anyone should deploy. We keep it because the
$1.6$--$1.8\times$ ratio is the paper's claim about $\rhoSym$, and that claim is
about usable memory, not about what the allocator will grant.

Two features of this table are worth stating precisely. The ratio is
\emph{flat} at $1.63$--$1.82\times$ across a twentyfold range of accuracy budgets,
which is what a budget-geometry argument predicts: both devices are compressing
the same side of the same model, so the constraint scales both answers together.
And the weight column is pinned at 8 bits everywhere. The solver is free to
choose any per-layer width and declines to spend on weights --- not because we
forbade it, but because the measured curves make the cache the better buy at
every budget shown (Table~\ref{tab:bitcost}).

\begin{table}[t]
\centering
\caption{Four configurations deployed at $L{=}4096$ on Device~A over four
windows of 1{,}024 tokens. $\Delta$PPL is relative perplexity against the
\texttt{nf4} row, so negative is better. ``full'' is full-precision weights with
an untouched cache: the clean control for the weight-precision effect, since no
cache quantisation sits on its path. The two cache-compressed rows are
\emph{simulated} --- our attention patch quantises and immediately dequantises ---
so their latency and memory are upper bounds rather than what a packed 4-bit
implementation would reach, and no claim in this paper rests on them.}
\label{tab:headtohead}
\begin{tabular}{@{}llrrrr@{}}
\toprule
 & config & $\Delta$PPL (\%) & MB & TTFT & TPOT \\
\midrule
0.6B & nf4     & $0.00$ & 2{,}166 & 450 & 32.2 \\
     & policy  & $\mathbf{-9.97}$ & 2{,}789 & 465 & 92.9 \\
     & both    & $+2.22$ & 2{,}167 & 516 & 93.2 \\
     & \textbf{full} & $\mathbf{-10.84}$ & 2{,}810 & \textbf{402} & \textbf{27.0} \\
\midrule
1.7B & nf4     & $0.00$ & 2{,}949 & 881 & \textbf{32.5} \\
     & policy  & $\mathbf{-4.68}$ & 4{,}942 & 859 & 100.8 \\
     & both    & $+0.08$ & 2{,}948 & 954 & 96.1 \\
     & \textbf{full} & $\mathbf{-4.68}$ & 4{,}998 & \textbf{794} & 35.3 \\
\bottomrule
\end{tabular}
\end{table}

\paragraph*{What the default actually costs, measured end to end}
Two natural ways to make this comparison are both wrong, and the errors are
instructive. Pricing the weight side from a cost curve at our own 4-bit
round-to-nearest makes the default look catastrophic ($36.5\%$); real \texttt{nf4}
is a far better quantiser than plain rounding, so that is a strawman. Comparing
\texttt{nf4} weights against full-precision weights that carry a \emph{simulated}
4-bit cache makes the default look three times faster; that measurement includes a
software dequantisation step on every layer and therefore describes our patch
rather than the default.

Table~\ref{tab:headtohead} carries the control that settles it: full-precision
weights with an untouched cache, so the only difference from the \texttt{nf4} row
is the weight precision. Three things follow.

First, on 0.6B at 4k the published practice is strictly dominated. It is $10.84\%$
worse in perplexity, $12\%$ slower to first token, $19\%$ slower per decode step,
and saves 644\,MB --- which the device does not need, since both configurations
fit. There is no dimension in which it wins.

Second, on 1.7B it becomes a real trade, because 3.4\,GB of fp16 weights is worth
compressing: \texttt{nf4} saves 2\,GB and decodes $9\%$ faster at a cost of
$4.68\%$ perplexity. That is a defensible engineering choice at 4k, and we do not
claim otherwise.

Third, compressing both sides is not the answer either: the row labelled
\emph{both} lands within $0.1$--$2.2\%$ of the default while the cache is small
enough at this length for the saving to be pointless.

The two cache-compressed rows are simulated: our attention patch quantises and
immediately dequantises, so it pays the arithmetic cost of compression without
storing packed bits. Their latency and memory therefore bound what a packed
implementation would cost, not what it would achieve, and no accuracy claim in
this paper rests on them.

The \emph{memory} claim, however, is not left modelled. We ran a prefill at 8k and
32k, captured every layer's key and value tensors, packed them to 4 bits two
values per byte with a per-32-group scale and minimum, and asked the allocator how
much it actually holds. At $L{=}8192$ the fp16 cache is 902\,MB and the packed form
is 282\,MB; at $L{=}32768$ they are 3{,}611\,MB and 1{,}129\,MB. Both are a
$3.2\times$ reduction --- below the ideal $4\times$ because a 5-bit budget covers
the nibbles plus one bit of scale, which is the arithmetic the capacity table
assumes. What this does not demonstrate is a decode step running on the packed
form; that needs a chunked dequantising attention kernel, which we did not write,
and the capacity argument needs only the resident footprint.

What makes the default the wrong choice regardless is the length at which it stops
fitting. The cache grows with $L$ and the weights do not, and on this laptop the
driver spills to system memory rather than raising an error, so the failure is
gradual: at 32k the 0.6B \texttt{nf4} configuration already peaks at 13.7\,GB
against 6.9\,GB of dedicated memory, and its decode time per token rises from
$32$\,ms to $1{,}076$\,ms as a result. A packed 4-bit cache would need
$1.8$\,GB where 16-bit needs $7.2$\,GB at that length. The honest summary is
conditional: \emph{at 4k the default is a modest trade on the larger model and a
dominated choice on the smaller one; past roughly 16k on a 6.9\,GB budget it is
not a trade at all, and cache compression is what makes the length reachable}.

\paragraph*{Is the policy an artefact of a weak weight quantiser?}
That is the obvious objection, and it separates from the ordering question of
Section~\ref{sec:results-bitcost}: even if a cache bit is cheaper, a good enough
weight quantiser might still be worth using. We tested it directly by rebuilding
the allocator's inputs with the weight side at the \emph{tuned} cost of
Table~\ref{tab:bitcost} ($6.97\%$ at 4 bits instead of $34.65\%$, a factor of
five) and re-solving. Almost nothing moves: at $\epsilon{=}0.01$ both devices
still choose 8-bit weights and the attainable context is unchanged to the token;
at $\epsilon{=}0.20$ the mean weight width is $7.98$ on the laptop and $8.00$ on the Jetson
with the default curves, against $7.96$ and $7.98$ with the tuned ones. The reason is geometric rather
than numerical: at these context lengths the cache is the dominant consumer of
bytes, so a bit moved to the cache buys far more context than the same bit buys
in accuracy. The policy is therefore robust to the strongest objection we can
raise against it, even though the underlying ordering is not
(Section~\ref{sec:results-bitcost}).

\subsection{The cache is not the fragile side}
\label{sec:results-bitcost}

Quantisation practice spends its bits on weights and leaves the cache in 16
bits, on the implicit assumption that the cache is the part that breaks first.
Table~\ref{tab:bitcost} measures both sides the same way --- quantise every
linear layer at a given width and record the perplexity cost, then quantise the
cache at the same width and record the cost --- and finds the opposite: at 4 bits
the cache costs $5.35\%$ and the weights cost $34.65\%$; at 3 bits the cache costs
$25.2\%$ while the weights collapse outright. Both intervals exclude zero, so the
ordering is not a small-sample accident.

We repeated the comparison on a second checkpoint, Qwen3-1.7B, because a single
model is the obvious objection to any claim of this shape. The ordering
reproduces and the ratio is close ($5.7\times$ against $6.5\times$), although
both costs fall in absolute terms ($11.33\%$ and $1.98\%$): the smaller model is
the more weight-sensitive one, which is consistent with weight quantisation
mattering less as models grow. The tuned margin also reproduces at about one
percentage point, so the two models agree on both the direction and the size of
the effect.

\begin{table}[t]
\centering
\caption{Relative perplexity increase when all weights, or the whole KV cache,
are quantised to the same width, on two model sizes. Corpus: 2{,}000 arXiv
abstracts, 48 windows of 1{,}024 tokens; intervals are 95\% bootstrap CIs over
windows. The ``tuned'' row strengthens the weight side as far as group size and
grid symmetry allow and is the honest worst case for our own claim. The
ordering --- cache cheaper than weights at the same width --- reproduces at
1.7B, but the margin under a tuned quantiser shrinks to about one percentage
point on both.}
\label{tab:bitcost}
\begin{tabular}{llll}
\toprule
bits & weights & KV cache & ratio \\
\midrule
\multicolumn{4}{l}{\emph{Qwen3-0.6B}}\\
8 & $-0.08\%$ & $0.00\%$ & --- \\
4 & $+34.65\%$ & $+5.35\%$ & $6.5\times$ \\
  & {\small[33.2, 36.1]} & {\small[4.9, 5.9]} & \\
4 (tuned) & $+6.97\%$ & $+5.35\%$ & $1.30\times$ \\
  & {\small[6.4, 7.5]} & {\small[4.9, 5.9]} & \\
3 & collapse & $+25.23\%$ & --- \\
2 & collapse & collapse & --- \\
\midrule
\multicolumn{4}{l}{\emph{Qwen3-1.7B}}\\
4 & $+11.33\%$ & $+1.98\%$ & $5.7\times$ \\
  & {\small[10.1, 12.5]} & {\small[1.4, 2.5]} & \\
4 (tuned) & $+3.21\%$ & $+1.98\%$ & $1.62\times$ \\
  & {\small[2.6, 3.8]} & {\small[1.4, 2.5]} & \\
\bottomrule
\end{tabular}
\end{table}

The tuned rows matter for how strongly this result can be stated. Our weight
quantiser is round-to-nearest, while GPTQ and AWQ extract substantially more from
the same width, so a reviewer is entitled to ask whether the ordering is an
artefact of a weak quantiser. Strengthening the weight side as far as group size
and grid symmetry allow brings its cost from $34.7\%$ down to $6.97\%$ --- close
to the cache's $5.35\%$, but not below it: the paired difference, measured on the
same windows, is $1.5$ percentage points with a 95\% interval of
$[-2.16, -0.84]$, which excludes zero but only just. We therefore claim the
weaker, safer statement: \emph{with quantisers of comparable strength the cache
is not the fragile side}, so treating it as untouchable while compressing weights
lacks a basis. The published GPTQ implementation, run at the same width
below, does not close that margin either --- it widens it --- but we still do not
claim that a cache bit is always cheaper. The
un-tuned $34.7\%$ figure is reported for completeness but should not be compared
against the cache number: it measures round-to-nearest, not the state of the art.

\begin{table}[t]
\centering
\small\caption{The same comparison under three quantiser pairs, each stronger than the
last, with calibration windows held out from evaluation. The cache is cheaper per
bit in every symmetric pair, by between $3.2\times$ and $14.3\times$. The single
row where it is not --- weights at group 16 against an ungrouped cache --- is an
artefact of giving one side finer granularity than the other; we report it
because we produced it ourselves and it is the natural way to get this comparison
wrong.}
\label{tab:strength}
\begin{tabular}{@{}llrrr@{}}
\toprule
 & & \multicolumn{2}{c}{rel.\ PPL (\%)} & ratio \\
\cmidrule(lr){3-4}
weights & cache & 0.6B & 1.7B & 0.6B / 1.7B \\
\midrule
RTN g128 & per-ch/tok & 34.71 & 11.45 & 6.2 / 6.2 \\
AWQ g128 & KIVI R128 & 16.45 & 9.61 & 3.2 / 4.5 \\
AWQ g16 & KIVI+g16 & 3.27 & 2.75 & 4.2 / 14.3 \\
\midrule
AWQ g16 & KIVI R128 & 3.27 & 2.75 & \textbf{0.6 / 1.3} \\
\multicolumn{5}{@{}l@{}}{\small\emph{asymmetric: weights grouped, cache not}}\\
\bottomrule
\end{tabular}
\end{table}

\paragraph*{Does the ordering survive stronger quantisers on both sides?}
The tuned row above strengthens only the weights, so it answers whether our
conclusion is a weak-weight-quantiser artefact. It does not answer the symmetric
question: if the cache is also compressed by a state-of-the-art method, does it
stay the cheaper side? We implemented both: activation-aware per-channel scaling
in the spirit of AWQ, with the per-layer exponent chosen on held-out calibration
windows, and a KIVI-style cache with a $128$-token full-precision residual window
and optional grouping. Table~\ref{tab:strength} reports the grid.

Two things follow. First, the ordering is stable under symmetric strengthening:
the cache is between $3.2\times$ and $14.3\times$ cheaper per bit across three
quantiser pairs and two models, and the ratio does not drift monotonically with
strength --- it depends on the pair, not on a single side. Second, the one
reversal in the table is instructive. Giving the weights group-16 while leaving
the cache at its published granularity makes the cache look $1.6\times$
\emph{more} expensive on 0.6B. That is not a property of the cache; it is an
artefact of unequal treatment, and we produced it ourselves before adding
matching grouping on the cache side. Any claim in this area should state the
granularity on both sides, and we now do.

\paragraph*{Both sides from published implementations}
Everything above uses at least one quantiser of ours, so the last check removes
both: the weights are quantised by the published \texttt{gptqmodel} (v4.2.5,
second-order error feedback, 4 bits, group 128) and the cache by KIVI's own
quantiser (their \texttt{utils\_quant}, per-channel keys with sequence grouping,
per-token values with head-dimension grouping, and a 128-token residual window)

The result is not marginal. At 0.6B the GPTQ weights cost $+11.39\%$ against a
cache cost of $+1.27\%$ --- a $9.0\times$ margin. At 1.7B the same pair gives
$+3.20\%$ against $+0.24\%$, a $13.3\times$ margin. For comparison, our own
KIVI-style adapter costs $+1.65\%$ and $+0.95\%$ on the two models, so their
grouping is better than ours, which is the expected direction and the reason to
use their code.

Two implementation notes belong in the record. KIVI ships no Qwen3 model file, so
the attention plumbing is ours in both cases; and we run their quantiser in its
simulation path rather than through their Triton kernels, which are not available
on Windows. The comparison is therefore of quantisation schemes at equal width,
not of kernel implementations, and the paper claims the former.

One cross-device detail is worth recording because it licenses comparing numbers
measured on the two machines: the fp16 reference NLL agrees to five decimals
($3.850999$ on Device~B against $3.850961$ on Device~A). Quantisation cost is a
property of the weights, not of the GPU, and this is the check that says so.

\paragraph*{Does the cache's cost accumulate with context?}
One might expect it to: the cache is quantised per token, so a longer context
means more quantised elements, and prior work reports error accumulation over
autoregressive decoding. We tested it by repeating the comparison at window
lengths from 512 to 4{,}096 tokens with the token budget held at 32k windows in
total. The cache's relative cost does not grow: it is $6.71\%$ at $L{=}512$ and
$6.85\%$ at $L{=}4{,}096$, a change of $0.14$ points, while the weight cost moves
by more ($32.6\%$ to $37.9\%$) with no trend. The ratio therefore stays between
$4.9$ and $6.3$ across the range. Two consequences: the policy of
Table~\ref{tab:capacity} does not need a context-length caveat, and the
error-accumulation result that motivates residual windows in other work does not
show up as a context-length effect in perplexity here, which is itself worth
noting.

The allocator built on these curves reaches a striking policy from the other
direction: it is free to choose any per-layer width, and at every accuracy budget
we tried it leaves the weights at 8 bits and spends the budget on the cache
(Table~\ref{tab:capacity}). Two caveats bound that result: our cache quantiser
is the simple per-channel/per-token scheme with no residual window, so a
stronger one would make the cache cheaper still and push the same way; and the
corpus is abstract-level, so only relative perplexity is interpretable.

\begin{figure}[t]
\centering
\includegraphics[width=0.75\columnwidth]{figures/fig_bitcost.pdf}
\caption{Perplexity cost of quantising all weights versus the whole cache. At 3
bits the weights collapse while the cache still costs only $25\%$; below that the
comparison stops being meaningful.}
\label{fig:bitcost}
\end{figure}

\subsection{Negative results}
\label{sec:results-negative}

Quantisation has no effect on latency here (Section~\ref{sec:results-nospeed}),
so any claim that a bit-allocation scheme accelerates edge inference must be
justified by a different mechanism. When the accuracy budget is loose relative to
the device budget, the optimum is the highest precision that fits and no
allocation policy matters; the problem only becomes interesting when capacity
binds. We report both regimes.

\subsection{Threats to validity}
Three confounds could make a measured number mean something other than it
appears to, and we address each in the protocol rather than the prose.

The 4-bit path is bitsandbytes nf4, whose dequantisation is not fused into the
matmul; a kernel-fused quantiser could show a speedup that we do not, so the
``no speedup'' result is a statement about this stack rather than about 4-bit
arithmetic in general. Background load moved a short-context prefill time by
$10\%$ between two consecutive runs in our own measurements, which is why
compared configurations are interleaved and every ratio in
Table~\ref{tab:nospeed} comes from paired runs. And the latency tables are
computed on an architecture-exact random-weight model for the cases where no real
checkpoint was available; we verified that this reproduces real-weight memory
within $1\%$ and naive-path prefill within $4\%$ before relying on it, and every
number in Tables~\ref{tab:stack} and~\ref{tab:nospeed} is from real weights.
Scope limitations --- two model sizes and two corpora for the accuracy curves, one
device actually run end to end --- are separate and listed in
Section~\ref{sec:limitations}.

% --------------------------------------------------------------------------- %
\section{Reproducibility}
\label{sec:repro}

Every number in this paper is produced by a script in the artifact, and every
table is generated from a measurement file rather than transcribed. A single
command rebuilds the figures, the summary CSV, the allocator inputs, and the
manuscript, and runs the checks below; it currently passes end to end.

\begin{itemize}
\item \emph{Static manuscript checks} for unbalanced environments, table column
      counts, math delimiters, undefined references and missing bibliography
      entries. These catch the mechanical errors that would otherwise cost a
      submission cycle.
\item \emph{A numeric audit} that extracts every factor and percentage claim from
      the manuscript text and compares the headline ones against the result
      files. We added it after a revision changed a table without changing the
      abstract that quoted it; the audit found four such stale numbers.
\item \emph{A paired bootstrap} over evaluation windows for the weight-versus-cache
      comparison, so the ordering is stated with an interval instead of a point.
\item \emph{A model-free reproduction} of the attention-dispatch finding
      (Section~\ref{sec:results-stack}) that needs no checkpoint: it constructs
      the tensors, runs both backends, and prints the two peak-memory figures.
\item \emph{Solver checks}: invariants plus a brute-force enumeration on small
      instances, which confirms the dynamic program is exact when the byte grid
      is fine enough rather than merely plausible.
\end{itemize}

We also record the execution stack with every measurement, including the package
indices that were reachable from our network, because
Section~\ref{sec:results-stack} shows that the stack can change attainable
context by $4\times$ --- an effect larger than any allocation decision in this
paper.

\paragraph*{A task-level check where the scheme decides}
Perplexity is an average; it cannot tell whether a model still finds something it
was told 11{,}000 tokens ago, which is what a compressed cache threatens. We ran a
passkey retrieval set --- five prompts, one hidden passkey each --- through five
cache schemes (Table~\ref{tab:passkey}).

\begin{table}[t]
\centering
\caption{Passkey retrieval over prompts of about 11k tokens, Qwen3-0.6B, five
prompts per row. ``grouped'' is KIVI's per-channel keys, per-token values and
128-token residual window; ``ungrouped'' is plain per-channel/per-token rounding.
At four bits the task does not separate the schemes; at two bits it separates them
completely.}
\label{tab:passkey}
\begin{tabular}{llr}
\toprule
cache & width & recovered \\
\midrule
none (fp16) & 16 & 5/5 \\
ungrouped & 4 & 5/5 \\
grouped (KIVI) & 4 & 5/5 \\
grouped (KIVI) & 2 & 5/5 \\
ungrouped & 2 & \textbf{0/5} \\
\bottomrule
\end{tabular}
\end{table}

Two things follow. The first is reassuring and is why we keep the 4-bit
configurations in the paper: at four bits, the scheme does not matter for this
capability, so the perplexity differences we measure at that width are differences
in fluency rather than in retrieval. The second is a warning that belongs in any
allocation policy: at two bits the choice of scheme decides everything. The same
budget that KIVI's grouping and residual window carry without loss destroys the
capability outright when the rounding is naive --- and note that the naive 2-bit
configuration is not obviously broken by perplexity alone in a short-context
evaluation. Compressing the cache is the right direction, but only with the
scheme that keeps the recent tokens exact.

\paragraph*{The ordering reproduces on a standard corpus}
The results above use 2{,}000 arXiv abstracts, which invites the objection that the
numbers are an artifact of the corpus. We therefore repeated the strength-matched
comparison on the WikiText-2 test set, which became reachable through a mirror
after the rest of the study was complete (Table~\ref{tab:wikitext}).

\begin{table}[t]
\centering
\caption{The strength-matched comparison repeated on the WikiText-2 test set, same
window protocol, same code, same calibration discipline. The ordering holds on both
models and the margin is \emph{larger} than on the abstract corpus, so the
conclusion does not depend on the corpus.}
\label{tab:wikitext}
\begin{tabular}{lrr}
\toprule
 & Qwen3-0.6B & Qwen3-1.7B \\
\midrule
weights, round-to-nearest & $+37.20\%$ & $+28.49\%$ \\
weights, AWQ-style (tuned) & $+17.52\%$ & $+6.33\%$ \\
cache, simple & $+5.25\%$ & $+1.30\%$ \\
cache, KIVI-style & $+5.34\%$ & $+2.42\%$ \\
cache, KIVI-style + grouping & $\mathbf{+1.61\%}$ & $\mathbf{+0.82\%}$ \\
\midrule
matched ratio (weights / cache) & $10.9\times$ & $7.7\times$ \\
\bottomrule
\end{tabular}
\end{table}

On WikiText-2 the matched weights cost $17.52\%$ against $1.61\%$ for the cache on
0.6B ($10.9\times$), and $6.33\%$ against $0.82\%$ on 1.7B ($7.7\times$). Both
margins are larger than the $3.2\times$--$14.3\times$ observed on the abstract
corpus, and the direction is identical, so the per-bit ordering is a property of
the models and the quantisers rather than of the text. The deployment and passkey
experiments remain on the abstract corpus so that they stay comparable with the
measurements taken before WikiText-2 was available; what changed is that the
central claim now has a standard-benchmark reproduction.

\section{Limitations}
\label{sec:limitations}

We list every limitation we know of, because each bounds a specific claim rather
than the paper as a whole.

\begin{enumerate}
\item \textbf{Model coverage.} The accuracy curves are measured on Qwen3-0.6B and
      Qwen3-1.7B. We tried to add a 4B point and could not do it on equal terms: a
      full-precision 4B reference does not fit in 8\,GB, so the weight side would
      have been bnb nf4 against a bnb int8 reference while the cache side used
      plain rounding, and those quantisers are not comparable. Under that
      mismatched setup the cache looks \emph{more} expensive ($+9.46\%$ against
      $+6.66\%$), which is why the ordering claim is scoped to the two models
      where the strength match holds. We do not know where the ordering reverses.

\item \textbf{The third-party check is split across devices.} The 0.6B official
      GPTQ run is on the edge module and the 1.7B run on the laptop, because the
      1.7B fp16 weights are 3.4\,GB against 4.3\,GB of usable memory on the module
      and error-feedback quantisation cannot compress weights before quantising
      them. We report the fp16 reference NLL agreeing to five decimals on both
      machines, which is what licenses the comparison. Calibration size also
      matters more than we expected: moving from 2 to 16 calibration texts made
      GPTQ \emph{worse} in our setup ($11.39\%\to13.0\%$ at 0.6B,
      $3.20\%\to5.96\%$ at 1.7B). The cache costs $1.27\%$ and $0.24\%$, so the
      ordering survives either way, but the absolute weight costs are
      calibration-sensitive and we report the range rather than one number.

\item \textbf{Where the task-level result applies.} On 11k-token passkey
      retrieval, every 4-bit scheme recovers the key, so at that width the
      perplexity differences we measure are differences in fluency rather than in
      retrieval. The schemes separate only at 2 bits, where KIVI's grouping and
      residual window still recover it 5/5 and ungrouped rounding recovers it 0/5.
      We therefore claim the task-level result only for the width at which we
      measured it, and the deployment comparison remains on the arXiv abstract
      corpus for comparability with the measurements taken before WikiText-2 was
      reachable. Absolute perplexities are not comparable to published numbers on
      either corpus; only relative increases are.

\item \textbf{Capacity beyond the measured range is analytic, and the edge column
      is an upper bound.} Table~\ref{tab:capacity} extrapolates from measured
      $\rho$ and a workspace curve; direct measurement contradicts it on the edge
      module, which runs 0.6B at 3{,}072 tokens and aborts its allocator at
      4{,}096 where the model predicts 67k. The
      column is retained because the ratio between devices is the claim about
      usable memory, but the text states plainly that the figure is not a deployable
      configuration. The workspace coefficient is also model- and
      stack-dependent, so it would not transfer to a different framework.

\item \textbf{The cache-compression measurement is partly simulated.} Our
      attention patch quantises and immediately dequantises, so its latency and
      peak memory bound what a packed implementation would cost rather than what it
      would achieve, and no accuracy claim rests on those rows. The memory claim is
      measured separately and directly --- packing the captured cache reduces its
      resident footprint $3.2\times$ at 8k and 32k --- but a decode step running on
      the packed form would need a chunked dequantising attention kernel, which we
      did not write. Relatedly, no published end-to-end allocation method is run as
      a baseline, since none of the systems we surveyed releases one; the
      comparison is against configurations, not against other allocators.

\item \textbf{No kernel engineering.} The gains reported here come from allocation
      and from repairing a dispatch failure, not from custom kernels. A fused
      quantised-attention kernel would change both the constants and, as
      Section~\ref{sec:results-stack} shows for the dispatch case, possibly the
      conclusions.
\end{enumerate}

\appendix
\section{Pre-registration and Two Falsified Hypotheses}
\label{app:prereg}

We wrote down what we expected to find before running the main experiments, and
we report the two hypotheses that failed. The point is not ceremony: the paper's
final claim is different from the one we started with, and a reader deserves to
know which of our statements were predictions and which were fitted to the data
after the fact.

\paragraph*{H1 (rejected): the optimal weight/KV split reverses between prefill-
and decode-bound regimes} The argument was that a discrete GPU and a unified-
memory module have different roofline crossovers $b^* = P_{\mathrm{eff}} /
B_{\mathrm{eff}}$, so the same workload lands in different regimes. We
pre-registered a prediction for the crossover on the second device and it was
wrong, in the opposite direction: we expected $b^*_{\mathrm{Jetson}} \in
[150, 280]$, and measurement gave $83$ against $124$ on the laptop
(4.51\,TFLOP/s over 54.1\,GB/s versus 22.36 over 161.9). The module is
\emph{compute-poor} rather than bandwidth-poor relative to the discrete GPU, so
the two crossovers are closer together than assumed and the regime window we
hoped to exploit is narrower, not wider. Three further observations killed the
hypothesis independently. Prefill processes the whole prompt as one batch, so for
any long prompt both devices are compute-bound and $b^*$ is irrelevant. The
decode time model is the same function of bytes on both devices, so an identical
error-versus-bytes frontier yields an identical choice. And the compute penalty
of low bit-widths shrinks with batch, so the mechanism that was supposed to drive
the asymmetry disappears exactly where it was needed.

\paragraph*{H2 (rejected): low bit-width trades compute for bandwidth} We
predicted $\eta_c(\text{nf4}) \le 0.85$, i.e.\ that 4-bit weights would slow
prefill. On synthetic weights the effect was undetectable; on the real checkpoint
nf4 is $1.21\times$ slower than fp16 at 1k context and $1.07\times$ at 16k
(Table~\ref{tab:nospeed}). The hypothesis was wrong in the opposite direction
from what we expected: low bit-width does not buy bandwidth at the cost of
compute, it costs both, because the dequantisation is unfused. This is why the
paper optimises capacity instead of latency.

\paragraph*{What we did not do} We did not re-measure with a different batch size
or a tighter error threshold to recover H1, and we did not change the objective
and re-run H2. Both are reported as they came out, and the negative-result
subsection states the regime in which bit allocation does not matter at all.

\section{Three Measurement Traps}
\label{app:trap}

Each of these produced plausible-looking but wrong numbers before it was caught,
and we report them because a reader re-running our harness would hit the same
traps.

\paragraph*{Timing outside inference mode} Every timed call must run under
\texttt{torch.inference\_mode}. Outside it, autograd retains activations for a
backward pass that never happens; the measured footprint inflates by roughly
$4\times$ and a spurious out-of-memory appears at 4k context for a 0.6B model
that fits 16k comfortably.

\paragraph*{Projecting the language-model head over the whole prompt} With a
151{,}936-token vocabulary the head over all positions is gigabytes, and it
scales linearly in $L$ --- larger than the weights and the cache combined at
4k. Real engines project one position. We verified the fix by a depth sweep: the
excess is constant in the number of layers (2{,}712\,MB at depth 1,
2{,}838\,MB at depth 8), which proves it is a single buffer rather than
per-layer accumulation.

\paragraph*{Silent spill} Windows WDDM backs an over-budget allocation with
system memory instead of raising. Throughput collapses by an order of magnitude
and nothing is logged. The harness therefore compares every peak against device
memory and marks the row; Table~\ref{tab:stack} shows such rows as
\emph{spilled} rather than reporting their timings.

\bibliographystyle{elsarticle-num}
\bibliography{refs}


\end{document}
